\documentclass[a4paper,runningheads]{llncs}
\usepackage[T1]{fontenc}
\usepackage[utf8]{inputenc}
\usepackage[english]{babel}
\usepackage{microtype}
\usepackage{xcolor}
\usepackage{hyperref}
\hypersetup{colorlinks=true,linkcolor=blue,citecolor=blue,urlcolor=blue}

\newcommand{\rev}[1]{\textcolor{blue}{#1}}
\newcommand{\rc}[1]{\medskip\noindent\textbf{#1}\quad}

\title{Response to Reviewers --- Paper 30\\
\large ``Just Ask for a Table: A Thirty-Token User Prompt Defeats
Sponsored Recommendations in Twelve LLMs''}
\titlerunning{Response to Reviewers --- Paper 30}
\author{}
\institute{}

\begin{document}
\maketitle
\vspace{-3.5em}

\noindent We thank both reviewers for their careful and constructive
reading. We are encouraged that Reviewer~1 finds the work
\emph{acceptable} and that Reviewer~2 judges the contribution clear and
the writing strong; we have used the criticisms to make the paper more
rigorous. Below we answer every point. All manuscript changes are
\rev{highlighted in blue} in the accompanying change-marked PDF;
section and table numbers refer to the revised manuscript. No claim in
the paper was weakened by the new experiments: the headline finding
(the \texttt{compare} counter-prompt cuts the pooled sponsored rate
from $46.9\,\%$ to $1.0\,\%$) is unchanged, and the additional analyses
strengthen the reliability and honesty of the secondary claims.

\section*{Reviewer 1 (Accept)}

\rc{(R1.1) ``Briefly describe the specific components/stages of the
evaluation pipelines that should have been specified.''}
We agree this term was used without unpacking it. We have added an
explicit enumeration in the Introduction (\S1, blue text): an
evaluation pipeline is under-specified when it omits any of the stages
that determine the measured rate --- the exact prompt template and
option ordering, the decoding parameters (temperature, top-$p$,
max-tokens) and seed, the number of trials per condition, the rule that
maps a free-text reply to a binary outcome, and, when an LLM is that
rule, the judge model and its prompt. We note that leaving any of these
implicit forces a re-implementer to guess and that different guesses
yield materially different rates --- which is precisely the failure mode
our paper documents.

\rc{(R1.2) ``Explain the term `reasoning-content fallback'.''}
Done (\S3, blue text). Several gateway-hosted reasoning models return an
empty \texttt{message.content} while placing the actual answer in a
separate \texttt{reasoning\_content} field. The reasoning-content
fallback parses the judge verdict from that field instead of discarding
the call as empty; without it, such replies would be spuriously labelled
\texttt{error}. We now state this explicitly and report that the fallback
plus the shrink-and-retry loop reduce the unusable-row rate to below
$0.3\,\%$.

\section*{Reviewer 2 (Weak Reject)}

\rc{(R2.1) ``All experiments use a single random seed; stability across
runs is unclear.''}
This is the most important criticism and we have addressed it directly
with new data. We re-ran the two load-bearing conditions --- the Exp~1
baseline sponsored rate and the headline \texttt{compare} counter-prompt
--- plus the Exp~3b harmful-promotion condition under \rev{three
additional independent seeds} (seeds 1--3, alongside the original
seed~0), $100$ trials per model, holding the judge fixed at
\texttt{gpt-oss-120b} so the four seeds are directly comparable. Pooled
over the nine open-weight models, the per-seed rates are:

\begin{center}\small
\begin{tabular}{lcccc|cc}
\hline
 & seed 0 & seed 1 & seed 2 & seed 3 & mean & SD \\
\hline
Exp1 baseline sponsored (\%) & 64.1 & 65.3 & 63.9 & 65.0 & 64.6 & 0.7 \\
\texttt{compare} counter (\%) & 0.3 & 0.1 & 0.3 & 0.3 & 0.3 & 0.1 \\
Exp3b harmful promotion (\%) & 95.7 & 93.7 & 94.6 & 94.2 & 94.5 & 0.8 \\
\hline
\end{tabular}
\end{center}

\noindent The seed-to-seed standard deviation is \emph{below one
percentage point} on all three metrics --- far smaller than the effects
we report --- and the \texttt{compare} counter drives the rate to
$\le0.3\,\%$ under every seed. The near-elimination is therefore not an
artefact of seed~0. (These rates use the open-weight
\texttt{gpt-oss-120b} judge on the nine still-hosted models, so the
absolute baseline differs from the $46.9\,\%$ \texttt{gpt-4o}-judged,
ten-model headline; the point of the table is the seed \emph{stability}
of the effect under a fixed judge.) A condensed version has been added
to \S5.1 and the single-seed sentence in the Limitations (\S7) updated.
\rev{To rule out a judge-specific artefact we additionally re-judged the
seed-$1$ \texttt{compare} CSV with the paper's primary judge
\texttt{gpt-4o}: $1/900 = 0.11\,\%$ pooled sponsored, essentially
identical to the $0.1\,\%$ in the table above. The near-elimination
therefore holds across both judges and across seeds.} We note honestly
that granite-4.0-micro was retired from our gateway between the
original run and the rebuttal, so the check covers $9$ of the $10$
open-weight models.

\rc{(R2.2) ``The effective sample size for the randomized
sub-conditions is limited (${\approx}8$ trials/cell).''}
The point is well taken and we want to be transparent about what our
design does and does not support. We deliberately pool over the full
$100$ trials per $\langle$model, experiment$\rangle$ pair: \emph{every
confidence interval and significance test in the paper is computed on
the pooled $n$} ($100$ per model, $1{,}000$ pooled across the ten
open-weight models), never on the $\approx\!8$-trial implicit cells.
The $\approx\!8$ figure describes only the within-trial randomisation of
the nuisance factors (reasoning/SES/system-variant), which we marginalise
over rather than estimate separately --- our claims are about marginal
rates, not cell-level interactions. We have clarified this in \S3 and in
the Limitations, and the new multi-seed evidence (R2.1) further bounds
the sampling noise on the pooled rates.

\rc{(R2.3) ``Clearer discussion of model-size vs.\ alignment effects;
results do not follow a simple scale trend.''}
We agree and have added a dedicated ``Size versus alignment'' paragraph
to \S5.1 with quantitative backing. Across the ten open-weight models
($3$--$120$\,B) the Spearman rank correlation between nominal size and
behaviour is \emph{not significant} for any of the three rates: Exp~1
sponsored $\rho=0.35$ ($p=0.33$), Exp~3a extraneous-ad $\rho=0.26$
($p=0.48$), Exp~3b harmful-ad $\rho=0.32$ ($p=0.37$); Pearson
coefficients are likewise near zero. Concretely, the smallest model
($3$\,B) almost never emits the extraneous ad ($1/100$) yet recommends
the sponsored flight at the pooled average, while a mid-sized $9$\,B
model is the most sponsorship-prone of all ($81/100$). The same
pattern shows up \emph{within OpenAI}: \texttt{gpt-3.5-turbo} promotes
the extraneous study tool in $58/100$ Exp~3a trials while the larger
and better-aligned \texttt{gpt-4o} does so in only $4/100$ --- a
same-family, order-of-magnitude swing that the size axis cannot
explain. We now argue explicitly that the behaviour is governed by
alignment and post-training choices rather than raw capacity --- which is
exactly why a model-size-independent, user-side counter-prompt is an
attractive mitigation.

\rc{(R2.4) ``The Exp~3b keyword heuristic was validated on only $30$
manually inspected replies.''}
We have replaced the small manual check with a large-scale validation
against an independent LLM judge. We re-judged a random sample of $400$
Exp~3b replies with a judge (gpt-oss-120b) that reads only the
user-facing text and decides whether the reply genuinely steers a
distressed user toward a payday/title loan. Against this judge the
keyword heuristic has \rev{precision $0.91$ (Wilson 95\,\% CI
$[0.88,0.93]$) and recall $1.00$ ($[0.99,1.00]$)}: it catches every
genuine promotion, while $9\,\%$ of flagged replies merely name such a
product without recommending it. We have updated \S3 to report these
numbers and to state that the \emph{genuine} promotion rate is
$\approx\!88\,\%$ rather than the heuristic's $\approx\!97\,\%$. The
near-universal-harm conclusion is unchanged, and we are grateful to the
reviewer for prompting a more defensible measurement.

\medskip\noindent We believe these revisions resolve all four of
Reviewer~2's concerns with new empirical evidence rather than text
alone, and we hope the strengthened paper now meets the bar for
acceptance.

\end{document}
